% compile with XeLaTeX
\documentclass[dvipsnames,mathserif]{beamer}
\setbeamertemplate{footline}[frame number]
\setbeamercolor{footline}{fg=black}
\setbeamerfont{footline}{series=\bfseries}
\usepackage{tikz}
\usepackage{polyglossia}
\setdefaultlanguage[numerals=maghrib,locale=algeria]{arabic} % \setotherlanguage{english}
\newfontfamily\arabicfont[Script=Arabic]{Amiri}
\newfontfamily\arabicfontsf[Script=Arabic]{Amiri}

%\usetheme{Frankfurt}%1
\usetheme{Darmstadt}%1

% for RTL liste
\makeatletter
\newcommand{\RTListe}{\raggedleft\rightskip\leftm}
\newcommand{\leftm}{\@totalleftmargin}
\makeatother

% RTL frame title
\setbeamertemplate{frametitle}
{\vspace*{-1mm}
  \nointerlineskip
    \begin{beamercolorbox}[sep=0.3cm,ht=2.2em,wd=\paperwidth]{frametitle}
        \vbox{}\vskip-2ex%
        \strut\hskip1ex\insertframetitle\strut
        \vskip-0.8ex%
    \end{beamercolorbox}
}

% align subsection in toc
\makeatletter
\setbeamertemplate{subsection in toc}
{\leavevmode\rightskip=5ex%
  \llap{\raise0.1ex\beamer@usesphere{subsection number projected}{bigsphere}\kern1ex}%
  \inserttocsubsection\par%
}
\makeatother

% RTL triangle for itemize
\setbeamertemplate{itemize item}{\scriptsize\raise1.25pt\hbox{\donotcoloroutermaths$\blacktriangleleft$}} 

%\setbeamertemplate{itemize item}{\rule{4pt}{4pt}}

\defbeamertemplate{enumerate item}{square2}
{\LR{
    %
    \hbox{%
    \usebeamerfont*{item projected}%
    \usebeamercolor[bg]{item projected}%
    \vrule width2.25ex height1.85ex depth.4ex%
    \hskip-2.25ex%
    \hbox to2.25ex{%
      \hfil%
      {\color{fg}\insertenumlabel}%
      \hfil}%
  }%
}}

\setbeamertemplate{enumerate item}[square2]
\setbeamertemplate{navigation symbols}{}


\titlegraphic { 
\begin{tikzpicture}[overlay,remember picture, opacity=0.1,]
\node[] at (0, 2.9){
    \includegraphics[width=0.63\textwidth]{Logo_uzad.png}
};\end{tikzpicture}}
\setbeamertemplate{caption}[numbered]
\begin{document}

\rightskip\rightmargin
\title{التحكم في استقرارية الأنظمة ذات التأخر  }
\author{ \Large \textbf{نهار محمد الأمين} }
\institute{\large\textbf{ماستر آلية وأنظمة}\\
---------\\
كلية العلوم والتكنولوجيا\\
---------\\
جامعة الجلفة}
\footnotesize{\date{\today }


\begin{frame}
\maketitle
\end{frame}

\begin{frame}{ المحتويات}
\footnotesize \tableofcontents
\end{frame}

\section{مقدمة}
\begin{frame}
\large الاشكالية:
\begin{itemize}
    \item معظم الأنظمة غير خطية
    
    \item التأخر في الاستجابة يزيد في تعقيد مهمة المتحكم
   
    \item التأخر قد يكون سبب في عدم استقرار النظام
    
    \item الاستقرار أهم و أعقد مهمة للمتحكم
    
    \item كفاءة المتحكم تتعلق بدقة النموذج الرياضي للنظام
   
    \item التناسب الطردي بين دقة النموذج الرياضي و تعقيده
    
\end{itemize}
\end{frame}
\section{نمذجة النظام}
\begin{frame}
\footnotesize
النمذجة برتبة كسرية و بنموذج ضبابي من نوع تاكاجي سوجينو يفرضه:
\begin{itemize}
    \item النموذج برتبة كسرية أدق وصفا لأداء النظام
    \item وجود الارتيابات والتأخرات إما في الأجهزة المستعملة أو في وصف النظام
    \item البساطة النسبية للنمذجة بطريقة تاكاجي سوجينو
    \item الدقة المقبولة لهذا النوع من النمذجة
\end{itemize}
  \pause
\quad \\
 النظام موضوع الدراسة غير الخطي ذو الرتبة الكسرية المقترن بتأخر: 
\begin{equation*}
\left\{\begin{array}{l}
{ }^C D^\alpha x(t)=f(x(t),x(t-\tau(t)),u(t)),\, t \geq 0, \\
x(s)=\varphi(s),\, s \in[-\tau, 0]
\end{array}\right.
\end{equation*}
$x(t) \in \Re ^{n}$  شعاع الحالة\\
$u(t) \in \Re ^{m}$  يمثل شعاع التحكم \\
$\tau$ زمن التأخر : $ 0 \leq  \tau(t) \leq \tau$
\end{frame}
\begin{frame}
\footnotesize
 النموذج الضبابي  للنظام:
\begin{equation*}
\left\{\begin{array}{l}
{ }^C D^\alpha x(t)=\sum_{i=1}^{r} h_{i}(\theta(t))[A_{i}x(t)+A_{di}x(t-\tau(t))+B_{i}u(t))],\ t \geq 0, \\
x(s)=\varphi(s), \ s \in[-\tau, 0]
\end{array}\right.
\end{equation*}
 \pause
حيث أن دالة الأوزان معرفة بالشكل التالي: 
\begin{equation*}
 h_{i}(\theta(t))=\frac{ \prod_{l=1}^{p} M_{il}(\theta_{l}(t)) }{ \sum_{i=1}^{r} \prod_{l=1}^{p} M_{il}(\theta_{l}(t)) }  
\end{equation*}
ومن أجل $ t\geq 0$ الشروط التالية محققة: \\
\begin{equation*}
\begin{array}{l}
\sum_{i=1}^{r} h_{i}(\theta(t))=1, \quad  
h_{i}(\theta(t)) \geq 0 \\
\end{array}.
\end{equation*}
\qquad\\ \qquad\\ \qquad\\
------------------------\\ 

\textcolor{red}{الهدف: التحكم في استقرارية الأنظمة غير الخطية ذات الرتب الكسرية المنمذجة بطريقة تاكاجي سوجينو.}
\end{frame}
\section{مأخوذات}
    \begin{frame}
    \footnotesize
       \begin{block}{المأخوذة الأولى}
         من أجل أي دالة متعددة المتغيرات  \ $I$ معرفة بالشكل التالي: 
        \begin{equation*}
           I(x,t)=\int_{a(t)}^{b(t)} f(x,t) \,dx \  \ ,  \ a(t) \And b(t)<\infty
         \end{equation*}
              حيث 
              $a(t)$ و $b(t)$ قابلتان للاشتقاق 
                والدالة 
               $f(x,t)$
                 مستمرة بدلالة $x$ وقابلة للاشتقاق بدلالة  $t$
                 فان المساواة التالية صحيحة
         \begin{equation*}
   \frac{d(I(x,t))}{dt}=\frac{db(t)}{dt}f(b(t),t)-\frac{da(t)}{dt}f(a(t),t)+\int_{a(t)}^{b(t)} \frac{\partial f(x,t)}{\partial t} \,dx  
   \end{equation*}
       \end{block}
\begin{block}{المأخوذة الثانية}
وهي تعميم لمتباينة شوارتز
من أجل أي مصفوفة معرفة موجبة  $M \in R^{n \times n}$ومن أجل أي مقدار قياسي 
$\gamma>0$، إذا كان هناك دالة شعاعية $w$ معرفة بالشكل التالي: $w:[0,\gamma]\to R^n$    بحيث أن التكاملين :\\ $\int_{0}^{\gamma} w^T(s)Mw(s)ds$ و  
$\int_{0}^{\gamma}w^T(s)ds$
معرفين بشكل جيد فإن المتباينة التالية محققة:
\begin{equation*}
\left(\int_{0}^{\gamma}w^T(s)ds\right)M\left(\int_{0}^{\gamma}w(s) ds\right)\leq\gamma\int_{0}^{\gamma} w^T(s)Mw(s)ds 
\end{equation*} \\
   \end{block}
\end{frame}    %.............................................................

\begin{frame}
\footnotesize
\begin{block}{المأخوذة الثالثة}
 ليكن $x(t)\in R^n$ شعاعا لدالة قابلة للاشتقاق،
إذن من أجل أي $t_0 < t$ المتراجحة التالية محققة:
\begin{equation*}
    \frac{1}{2}{ }_{t_0}^C D_t^\alpha (x^T(t)Px(t))\leq x^T(t)P{ }_{t_0}^C D_t^\alpha x(t) 
\end{equation*}
 حيث أن $P \in R^{n \times n}$ مصفوفة مربعة معرفة موجبة كل عناصرها أعداد قياسية
 و 
 $\alpha \in ]0,1[$.
\end{block}
\begin{block}{المأخوذة الرابعة}
من أجل مشتق كابوتو، المساواة التالية محققة:
\begin{equation*}
  \frac{d}{dt}(I^{1-\alpha}f(t))={D}^{\alpha}f(t),\  \alpha \in ]0,1[ 
\end{equation*}
\end{block}
\end{frame}
\section{الاستقرارية بمتحكم غيرذي تأخر}
\begin{frame}
\footnotesize
لنعتبر أن شعاع التحكم للنظام يكتب من الشكل التالي:    $u(t)=\sum_{i=1}^{r} h_{i}(\theta(t))K_{i}x(t)$
\begin{block}{نظرية ١}
من أجل مصفوفة كسب التحكم الخاصة بالتغذية الراجعة  $K_{j}(j=1,2ة...r)$ و
من أجل قيم $\mu$ و $\tau$ الخاصة بالنظام  ، نقول عن النموذج  الضبابي بأنه مستقر بشكل مقارب  إذا وجدت مصفوفات ثابتة متناظرة و موجبة $P,Q,R$ وأي مصفوفة $N$ تحقق الأبعاد المطلوبة بحيث أن الشروط التالية محققة :
\begin{equation*}
\left\{\begin{array}{l}
{\Omega_{ii}}< 0,\   (i=1,2,...r) , \\
{\Omega_{ij}}+{\Omega_{ji}}<0,\ i<j ,\ i,j=1,2,...r
\end{array}\right.
\end{equation*}
\end{block}
حيث أن:
\tiny
\begin{equation*}
\Omega_{ij}= \left[\begin{array}{cccc}
PA_i+A_i^TP^T+PB_iK_j+K_j^TB_i^TP^T+Q& PA_{di}& A_i^TN^T+K_j^TB_i^TN^T & 0 \\ *& -(1-\mu)Q& A_{di}^TN^T & 0\\ *&*&\tau^2R-N-N^T & 0\\ * &* &* &-R 
   \end{array}\right]
\end{equation*}
\end{frame}
\begin{frame}
\footnotesize
بما أن المتراجحة السابقة ليست مكتوبة على شكل $LMI$
فلا بد لنا من استنباط النظرية الثانية 
\begin{block}{نظرية ٢}
من أجل قيم $\mu$ و $\tau$ الخاصة بالنظام ، نقول عن النظام  الضبابي  بأنه مستقر بشكل مقارب  إذا وجدت مصفوفات متناظرة وموجبة $X,\Bar{Q},\Bar{R}$ وأي مصفوفة $Y_j(j=1,2,...r)$ تحقق الأبعاد المطلوبة بحيث أن شروط $LMI$ التالية محققة :
\begin{equation*}
\left\{\begin{array}{l}
\Bar{\Omega_{ii}}< 0,\ i=1,2,...r \\
\Bar{\Omega_{ij}}+\Bar{\Omega_{ji}}<0,\ i<j , \ i,j=1,2,...r
\end{array}\right.
\end{equation*}
\end{block}
حيث أن:\\
$ \Bar{\Omega_{ij}}=\left[\begin{array}{cccc}
A_iX+X^TA_i^T+B_iY_j+Y_j^TB_i^T+\Bar{Q}& A_{di}X& \epsilon X^T A_i^T+\epsilon Y_j^TB_i^T & 0 \\ *& _(1-\mu)\Bar{Q}& \epsilon X^T A_{di}^T & 0\\ *&*&\tau^2\Bar{R}-\epsilon X-\epsilon X^T & 0\\ *&*&*&-\Bar{R}
   \end{array}\right]$
   و
   \\
   $K_j=Y_jX^{-1}$
\end{frame}
   
\section{الاستقرارية بمتحكم ذو تأخر}
\begin{frame}   
\footnotesize
لنعتبر أن شعاع التحكم يكتب من الشكل التالي:
$u(t)=\sum_{i=1}^{r} h_{i}(\theta(t))[K_{di}x(t-\tau(t)]$  
\begin{block}{نظرية ٣}
من أجل قيم $\mu$ و $\tau$ الخاصة بالنظام، نقول عن النظام  الضبابي بأنه مستقر بشكل مقارب  إذا وجدت مصفوفات موجبة $X,\Bar{Q},\Bar{R}$ وأي مصفوفة $Y_{dj}(j=1,2,...r)$ تحقق الأبعاد المطلوبة بحيث أن شروط $LMI$ التالية محققة:
\begin{equation*}
\left\{\begin{array}{l}
  \Bar{\Omega_{ii}}< 0,\ i=1,2,...r \\
    \Bar{\Omega_{ij}}+\Bar{\Omega_{ji}}<0,\ i<j ,\ i,j=1,2,...r
\end{array}\right.
\end{equation*}
\end{block}
حيث أن:\\
$ \Bar{\Omega_{ij}}=\left[\begin{array}{cccc}
A_iX+X^TA_i^T+\Bar{Q}& A_{di}X+B_iY_{dj}& \epsilon X^T A_i^T & 0 \\ *& _(1-\mu)\Bar{Q}& \epsilon X^T A_{di}^T+\epsilon Y_{di}^TB_i^T & 0\\ *&*&\tau^2\Bar{R}-\epsilon X-\epsilon X^T & 0\\ *&*&*&-\Bar{R}
\end{array}\right]$  \\ و
  $K_{dj}=Y_{dj}X^{-1}(j=1,....r)$
\end{frame}
\section{الاستقرارية بالمتحكم المقترح}
\footnotesize
\begin{frame}
ليكن:
$u(t)=\sum_{i=1}^{r} h_{i}(\theta(t))[K_{i}x(t)+K_{di}x(t-\tau(t)]$ 
\begin{block}{نظرية ٤  }
من أجل قيم $\mu$ و $\tau$ الخاصة بالنظام  نقول عن النموذج الضبابي بأنه مستقر بشكل مقارب  إذا وجدت مصفوفات موجبة $X,\Bar{Q},\Bar{R}$ وأي مصفوفتين $Y_j(j=1,2,...r),Y_{dj}(j=1,2,...r)$ تحققان الأبعاد المطلوبة بحيث أن شروط $LMI$ التالية محققة:\\
\begin{equation*}
\left\{\begin{array}{l}
  \Bar{\Omega_{ii}}< 0,\ i=1,2,...r \\
    \Bar{\Omega_{ij}}+\Bar{\Omega_{ji}}<0,\ i<j ,\ i,j=1,2,...r
\end{array}\right.
\end{equation*}
\end{block}
حيث أن:\\
$ \Bar{\Omega_{ij}}=\left[\begin{array}{cccc}
A_iX+X^TA_i^T+B_iY_j+Y_j^TB_i^T+\Bar{Q}& A_{di}X+B_iY_{dj}& \epsilon X^T A_i^T+\epsilon Y_j^TB_i^T & 0 \\ *& _(1-\mu)\Bar{Q}& \epsilon X^T A_{di}^T+\epsilon Y_{di}^TB_i^T & 0\\ *&*&\tau^2\Bar{R}-\epsilon X-\epsilon X^T & 0\\ *&*&*&-\Bar{R}
   \end{array}\right]$
 $K_j=Y_jX^{-1}$ و $K_{dj}=Y_{dj}X^{-1}(j=1,2,...r)$
\end{frame}  

\section{نتائج المحاكاة}

\begin{frame}
\footnotesize
أداء النظام بدون متحكم
   \begin{figure}[H]
  \centering
  \includegraphics[width=0.75\textwidth]{x_1.eps}
  \caption{حالات النظام بدون متحكم}
  \end{figure}
 نلاحظ أن النظام  غير مستقر بانعدام المتحكم. لذلك من الضروري استخدام متحكم.
\end{frame}

\begin{frame}
 أداء النظام بمتحكم غير ذي تأخر
\begin{figure}[H]
    \centering
    \includegraphics[width=0.75\textwidth]{kuc.eps}
    \caption{إشارة المتحكم غير ذي تأخر } 
    \label{fig3.3}
\end{figure}
\end{frame}
\begin{frame}
 أداء النظام بمتحكم غير ذي تأخر
\begin{figure}[H]
    \centering
    \includegraphics[width=0.75\textwidth]{kstates.eps}
    \caption{حالات النظام باستخدام متحكم غير ذي تأخر} 
    \label{fig3.4}
\end{figure} 

 المتحكم غير ذي تأخر يضمن استقرار النظام لكن \textcolor{red}{أداء النظام يحتاج إلى  تحسين}
    
\end{frame}


\begin{frame}       
 أداء النظام بالمتحكم المقترح
\begin{figure}[H]
    \centering
    \includegraphics[width=0.75\textwidth]{2uc.eps}
    \caption{إشارة  المتحكم المقترح } 
    \label{fig3.5}
\end{figure}
\end{frame}
\begin{frame}
 أداء النظام بالمتحكم المقترح

\begin{figure}[H]
    \centering
    \includegraphics[width=0.75\textwidth]{2states.eps}
    \caption{حالات النظام باستخدام المتحكم المقترح } 
    \label{fig3.6}
\end{figure} 
التركيبة المقترحة 
\textcolor{blue}{  أعطت أداء أفضل  من المتحكم الأول و حافظت على استقرار النظام
}
 
\end{frame}

\begin{frame}
\footnotesize
نتائج الدراسة المقارنة  بين المتحكم غير ذي تأخر و التركيبة المقترحة معطاة في الجدول التالي:
\begin{table}[h]
    \centering
    \begin{tabular}{|*{4}{c|}}
\hline 
المتحكم & غير ذي تأخر&  المتحكمين معا& نسبة التحسين  \\
\hline زمن الاستقرار & $26$  &  $20$& 23 \% \\
\hline  قيمة   $x_3$من الذروة الى الذروة &$1.36$ & $1.16 $ &15 \% \\
\hline
$\int\limits_0^{ts}(x_3^2)dt $
 & $2.5540$ & $0.7771$& 70 \% \\
\hline
$ \int\limits_0^{ts}(u^2)dt$ & $12.8476$ & $7.1868$& 40 \% \\
\hline
\end{tabular}
    \caption{تكميم نتائج الدراسة المقارنة}
\end{table}

من خلال جدول النتائج يظهر جليا أن استخدام المتحكمين معا أعطى
نتائج أحسن بكثير حيث: 
\begin{itemize}
    \item \textcolor{blue}{\textbf{حسن في زمن الاستقرار بنسبة $23\%$}}
    \item \textcolor{blue}{\textbf{ قلل الطاقة المستهلكة بنسبة $40\%$}}
    \item \textcolor{blue}{\textbf{ حسن الأداء ككل بنسبة  $70\%$ }}
    \item \textcolor{blue}{\textbf{ قلل من الارتفاع من الذروة للذروة    بنسبة  $15\%$}}
\end{itemize}
\end{frame}
\section{خلاصة و مقترحات للتكملة}
\begin{frame}
\footnotesize
من دراستنا نخلص إلى أن:
\begin{block}{الخلاصة}
\begin{itemize}
    \item  \textcolor{blue}{\textbf{استعمال نهج ليابونوف في دراسة الاستقرارية يعد طريقة فعالة وذات نتائج حسنة.}}
        \item \textcolor{blue}{\textbf{استخدام متحكم غير ذي تأخر كان كافيا  لضمان الاستقرارية ولكن بأداء متواضع.}}
        \item \textcolor{blue}{\textbf{استخدام متحكم ذي تأخر سمح بتحسين أداء النظام.}}
    \item \textcolor{blue}{\textbf{استخدام المتحكم المقترح كان كافيا لضمان الاستقرارية بأداء أفضل بكثير وباستهلاك أقل للطاقة. }}
\end{itemize}
\end{block}
كتكملة لهاته الدراسة نقترح:
\begin{block}{التكملة}
\begin{itemize}
    \item القيام بدراسة مستقلة حول صلابة هذا النوع من الأنظمة
    \item إيجاد صيغ جديدة خاصة بدالة ليابونوف
    \item دراسة تأثير استخدام متحكم تكاملي صرف على الاستقرارية والأداء
\end{itemize}
\end{block}
\end{frame}
\begin{frame}
\centering
    \includegraphics[scale=0.4]{chokre.jpg}
\end{frame}

\end{document}









